Dizemos que um inteiro positivo livre de quadrados \(n\) é \emph{quase primo} se
\[
n \mid x^{d_1} + x^{d_2} + \ldots + x^{d_k} - kx
\]
para todos os inteiros \(x\), onde \(1 = d_1 < d_2 < \ldots < d_k = n\) são todos os divisores positivos de \(n\). Suponha que \(r\) seja um primo de Fermat (ou seja, é um primo da forma \(2^{2^m} + 1\) para um inteiro \(m \ge 0\)), \(p\) é um divisor primo de um inteiro quase primo \(n\), e \(p \equiv 1 \pmod{r}\). Mostre que, com a notação acima, \(d_i \equiv 1 \pmod{r}\) para todo \(1 \le i \le k\).

(Um inteiro \(n\) é chamado de \emph{livre de quadrados} se não é divisível por \(d^2\) para qualquer inteiro \(d > 1\).)
